\section{Background and Motivation}

The anticipated arrival of large-scale quantum computers threatens the public-key
cryptography in use today. In response, the National Institute of Standards and
Technology (NIST) has been standardising post-quantum schemes, and has run a dedicated
process for additional digital signatures to diversify the mathematical assumptions
underlying signature schemes. \textbf{MQOM} (``MQ on my Mind'')~\cite{mqomspec} is one
such candidate. Its security reduces to the hardness of the \emph{Multivariate
Quadratic} (MQ) problem, and it is constructed using the
Multi-Party-Computation-in-the-Head (MPCitH) paradigm, more precisely its
Threshold-Computation-in-the-Head (TCitH) formulation. MQOM advanced to the third round
of the NIST additional-signatures process; this project studies the second-round
specification, \textbf{MQOM~v2.1}.

A cryptographic scheme can be mathematically sound yet still be broken on a real device.
\emph{Side-channel attacks} exploit physical quantities --- most commonly the power
consumption or electromagnetic emanation of a device --- that depend on the secret data
being processed. \emph{Correlation Power Analysis} (CPA)~\cite{brier2004cpa} is a
canonical side-channel technique: the attacker predicts, for each hypothesis about a
small part of the secret, the value of a secret-dependent intermediate, maps that
prediction to a leakage model (such as the Hamming weight), and correlates the
prediction against measured leakage. The hypothesis whose prediction correlates best is
taken as the recovered secret.

While MQOM's mathematical security has received considerable attention, its resistance
to side-channel analysis is far less studied. This project asks a focused question:
\emph{can the permanent secret key $x$ of MQOM~v2.1 be recovered by a CPA-based attack
exploiting the power leakage of secret-dependent intermediate values?}

\section{Problem Addressed}

The goal of this project is the complete recovery of the permanent secret witness
$x = (x_1, \dots, x_n)$ of MQOM~v2.1 through a CPA-based technique, with the success
criterion $\hat{x} = x$ (exact recovery of the full vector), and to establish the
mathematical conditions under which such recovery is possible. Where a given attack
point does not yield full recovery, the objective is to characterise rigorously how
much of the key it recovers, under which leakage model, and why --- and, failing a
clean recovery, to identify the highest-potential paths supported by our own
simulation evidence.

\section{Research Constraints}

This is a \textbf{computer-only} study. We deliberately and permanently exclude: physical
hardware and target devices; power or electromagnetic measurement equipment; real power
or EM traces; physical side-channel experiments; the official MQOM reference
implementation; and any compilation, execution or testing of that implementation. We
rely instead on independent mathematical reimplementations of the relevant operations,
models of MQOM intermediate values derived from the specification, synthetic leakage
datasets generated under explicitly stated hypotheses, and standard statistical
analysis. Consequently, every result in this report is a statement about mathematical
models and synthetic data; none is a measurement of, or a demonstrated vulnerability in,
a physical MQOM implementation.

\section{Objectives and Contributions}

The concrete objectives of the project are:
\begin{itemize}
\item to reassess MQOM~v2.1's candidate leakage locations under the computer-only
constraint, keeping only those that expose the secret to known, varying public data as
required by CPA;
\item to construct independent, specification-faithful Python models of the surviving
attack points and of the finite-field arithmetic they use;
\item to formulate explicit hypothetical leakage models and the matching CPA variants,
and to determine --- analytically and by simulation --- whether each is
\emph{identifiable}, i.e.\ whether the correct secret can be distinguished from all
competing hypotheses;
\item to demonstrate recovery of the secret (in whole or in characterised part) under
clearly stated assumptions, with negative controls that guard against
success-by-construction; and
\item to evaluate robustness to leakage-model mismatch and noise, and to state honestly
the limits of what simulation can establish.
\end{itemize}

\section{Organisation of the Report}

Chapter~\ref{ch:literature} reviews the relevant literature, carefully separating
physical side-channel analysis from mathematical cryptanalysis and fault attacks.
Chapter~\ref{ch:new} presents the proposed work: the reassessed attack surface, the
decomposition of the secret, the leakage-model taxonomy, the CPA formulations, and the
simulation methodology. Chapter~\ref{ch:result} reports the simulation results obtained
so far. Chapter~\ref{ch:conclude} draws conclusions and sets out future work.
